% !TEX program = lualatex
% Generated by scripts/build_module_practice.py - do not edit by hand.
\DocumentMetadata{
  pdfstandard=UA-2,
  pdfversion=2.0,
  lang=en-US,
  tagging=on
}
\documentclass[11pt]{article}
\usepackage[letterpaper, margin=0.8in, top=0.6in, bottom=0.6in]{geometry}
\usepackage{fontspec}
\setmainfont{Fira Sans}
\usepackage{amsmath}
\usepackage{amssymb}
% Fira Sans has no U+2032; an unmapped glyph fails PDF/UA-2 (.notdef)
\usepackage{newunicodechar}
\newunicodechar{′}{\ensuremath{'}}
\usepackage{booktabs}
\usepackage{longtable}
\usepackage{array}
\usepackage{calc}
\usepackage{etoolbox}
\usepackage{xcolor}
\usepackage{graphicx}
\definecolor{accent}{HTML}{BF5700}
% pandoc emits \tightlist on compact lists
\providecommand{\tightlist}{\setlength{\itemsep}{0pt}\setlength{\parskip}{0pt}}
\usepackage{titlesec}
\titleformat{\section}{\large\bfseries\color{accent}}{}{0pt}{}
\titlespacing*{\section}{0pt}{14pt}{4pt}
\titleformat{\subsection}{\normalsize\bfseries}{}{0pt}{}
\titlespacing*{\subsection}{0pt}{10pt}{2pt}
\usepackage{hyperref}
\hypersetup{pdftitle={ECON 201 Solutions: Firms as Price Setters}, pdfauthor={Div Bhagia}, hidelinks}
\pagestyle{plain}
\setlength{\parindent}{0pt}
\setlength{\parskip}{4pt}
\begin{document}
{\LARGE\bfseries\color{accent} Solutions: Firms as Price Setters}\hfill{\large ECON 201}
\vspace{4pt}


\section{Lecture 6: What does it cost to make a car? Scale and the cost of production}\label{lecture-6-what-does-it-cost-to-make-a-car-scale-and-the-cost-of-production}

\subsection{1. Denise's bakery}\label{denises-bakery}

\textbf{(a)} The fixed cost is \$300 and the variable cost is \$2 per loaf, so \(C(Q) = 300 + 2Q\).

\textbf{(b)}

\begin{longtable}[]{@{}
  >{\raggedright\arraybackslash}p{(\linewidth - 6\tabcolsep) * \real{0.2500}}
  >{\raggedleft\arraybackslash}p{(\linewidth - 6\tabcolsep) * \real{0.2500}}
  >{\raggedleft\arraybackslash}p{(\linewidth - 6\tabcolsep) * \real{0.2500}}
  >{\raggedleft\arraybackslash}p{(\linewidth - 6\tabcolsep) * \real{0.2500}}@{}}
\toprule\noalign{}
\begin{minipage}[b]{\linewidth}\raggedright
Loaves per day, \(Q\)
\end{minipage} & \begin{minipage}[b]{\linewidth}\raggedleft
50
\end{minipage} & \begin{minipage}[b]{\linewidth}\raggedleft
100
\end{minipage} & \begin{minipage}[b]{\linewidth}\raggedleft
300
\end{minipage} \\
\midrule\noalign{}
\endhead
\bottomrule\noalign{}
\endlastfoot
Total cost, \(C(Q)\) & 300 + 100 = 400 & 300 + 200 = 500 & 300 + 600 = 900 \\
Average cost, \(C(Q)/Q\) & 400 / 50 = 8 & 500 / 100 = 5 & 900 / 300 = 3 \\
Marginal cost & 2 & 2 & 2 \\
\end{longtable}

Marginal cost is the cost of one more loaf, \$2, at every output. Average cost is \(2 + 300 / Q\): each loaf carries its own \$2 plus a share of the \$300, and that share shrinks as more loaves are baked.

\textbf{(c)} Average cost is \(2 + 300 / Q\). Setting \(2 + 300 / Q = 3\) gives \(300 / Q = 1\), so \textbf{\(Q = 300\) loaves a day}. The table in part (b) confirms it.

\textbf{(d)} - \textbf{Total cost} rises by \$300 at every output: \(C(Q) = 600 + 2Q\).
- \textbf{Average cost} rises, by \(300 / Q\): at 100 loaves it goes from \$5 to \$8.
- \textbf{Marginal cost does not change.} One more loaf still costs \$2 in flour and labor. The rent is paid whether that loaf is baked or not.

\subsection{2. Fixed or variable, and for how long}\label{fixed-or-variable-and-for-how-long}

\textbf{(a)} \textbf{Fixed:} the rent and the website fee. Both are the same whether the bakery bakes one loaf or a thousand.

\textbf{Variable:} flour, yeast, and butter; hourly staff; and electricity for the ovens. Each of these rises with the number of loaves baked.

\textbf{(b)} Beyond 300 loaves each extra loaf needs overtime labor, so \textbf{marginal cost rises above \$2}. It is a short-run cost because it comes from holding the ovens fixed: in the short run the stock of equipment cannot be changed. In the long run Denise can buy another oven, and the extra loaves would again cost \$2 each.

\subsection{3. Lena's screen-printing shop}\label{lenas-screen-printing-shop}

\textbf{(a)} First the total costs:
\[C(10) = 90 + \frac{10^2}{10} = 100, \qquad C(15) = 90 + \frac{15^2}{10} = 112.5, \qquad C(20) = 90 + \frac{20^2}{10} = 130\]

From 10 to 15 shirts:
\[\text{MC} = \frac{\Delta C}{\Delta Q} = \frac{C(15) - C(10)}{15 - 10} = \frac{112.5 - 100}{5} = \$2.50 \text{ per shirt}\]

From 15 to 20 shirts:
\[\text{MC} = \frac{C(20) - C(15)}{20 - 15} = \frac{130 - 112.5}{5} = \$3.50 \text{ per shirt}\]

\textbf{(b)} \textbf{No.} Marginal cost rises with output, from \$2.50 a shirt to \$3.50. With a fixed press and a fixed number of hours in the day, printing more means squeezing more out of the same equipment: overtime pay, a tired crew, more misprints. That is the short-run case from the lecture, where marginal cost rises while the equipment is fixed.

\subsection{4. Returns to scale}\label{returns-to-scale}

\textbf{(a)} - The brewery more than doubles its output: \textbf{increasing returns to scale}, also called economies of scale.
- The print shop exactly doubles it: \textbf{constant returns to scale}.
- The consulting firm less than doubles it: \textbf{decreasing returns to scale}, also called diseconomies of scale.

\textbf{(b)} \textbf{Yes, if it has fixed costs.} Constant returns mean the variable cost per poster stays the same, but a fixed cost such as the lease on the press is spread over more posters as output rises. With \(C(Q) = F + cQ\), average cost is \(c + F / Q\), which falls as \(Q\) rises even though \(c\) does not change.

\subsection{5. Multiple choice}\label{multiple-choice}

\medskip

\textbf{Q1.} \textbf{(b).} Marginal cost is the increase in total cost from one more unit, which is the slope of the cost function: \textbf{\$2}. Option (c) is the average cost.

\medskip

\textbf{Q2.} \textbf{(a), (b) and (d).} At \(Q = 0\) total cost is 0, so there are no fixed costs. The slope of the cost function is 2, so marginal cost is 2. Average cost is \(2Q / Q = 2\) at every output, so it does not fall, and it equals marginal cost everywhere. \textbf{(a), (b), and (d)} are correct.

\medskip

\textbf{Q3.} \textbf{(a), (c) and (d).} Decreasing returns mean doubling inputs \emph{less than} doubles output, so (b) is wrong. The other three follow from the definitions: with increasing returns the firm needs a less than proportional increase in inputs to raise output, so its cost per unit falls. \textbf{(a), (c), and (d)} are correct.

\medskip

\textbf{Q4.} \textbf{(b).} Network economies of scale are on the demand side: the product is worth more to each user because other users are connected to it. That is \textbf{(b)}. Option (a) is bargaining power, (c) is a fixed cost spread over more units, and (d) is an economy of scale in production.

\medskip

\textbf{Q5.} \textbf{(c).} Average cost is the constant cost per car plus a share of the fixed cost, so it is always above marginal cost. The share shrinks as more cars are produced, so average cost falls toward marginal cost without reaching it. \textbf{(c)} is correct.

\section{Lecture 7: Who can get away with raising prices? Demand and elasticity}\label{lecture-7-who-can-get-away-with-raising-prices-demand-and-elasticity}

\subsection{1. Nadia's kayak rentals}\label{nadias-kayak-rentals}

\textbf{(a)} \[Q = 60 - 2 \times 20 = 20\]

\textbf{(b)} At \$21:
\[Q = 60 - 2 \times 21 = 18\]

The percentage changes:
\[\begin{aligned}
\%\text{ change in price} &= \frac{21 - 20}{20} \times 100 = 5\% \\[0.3em]
\%\text{ change in demand} &= \frac{18 - 20}{20} \times 100 = -10\%
\end{aligned}\]

The elasticity:
\[\varepsilon = -\frac{-10}{5} = 2\]

\textbf{Demand is elastic at \$20}, because 2 is greater than 1.

\textbf{(c)} From \$25 to \$26, demand falls from \(60 - 2 \times 25 = 10\) to \(60 - 2 \times 26 = 8\) kayaks:
\[\begin{aligned}
\%\text{ change in price} &= \frac{26 - 25}{25} \times 100 = 4\% \\[0.3em]
\%\text{ change in demand} &= \frac{8 - 10}{10} \times 100 = -20\%
\end{aligned}\]
\[\varepsilon = -\frac{-20}{4} = 5\]

From \$10 to \$11, demand falls from \(60 - 2 \times 10 = 40\) to \(60 - 2 \times 11 = 38\) kayaks:
\[\begin{aligned}
\%\text{ change in price} &= \frac{11 - 10}{10} \times 100 = 10\% \\[0.3em]
\%\text{ change in demand} &= \frac{38 - 40}{40} \times 100 = -5\%
\end{aligned}\]
\[\varepsilon = -\frac{-5}{10} = 0.5\]

The formula \(\varepsilon = -\dfrac{P}{Q} \times \dfrac{\Delta Q}{\Delta P}\), with \(\Delta Q / \Delta P = -2\), gives the same answers: \(-\tfrac{25}{10} \times (-2) = 5\) and \(-\tfrac{10}{40} \times (-2) = 0.5\).

\textbf{Demand is elastic at \$25 and inelastic at \$10.}

\textbf{(d)} Demand falls from \(60 - 0.02 \times 2{,}000 = 20\) to \(60 - 0.02 \times 2{,}100 = 18\) kayaks:
\[\begin{aligned}
\%\text{ change in price} &= \frac{2{,}100 - 2{,}000}{2{,}000} \times 100 = 5\% \\[0.3em]
\%\text{ change in demand} &= \frac{18 - 20}{20} \times 100 = -10\%
\end{aligned}\]
\[\varepsilon = -\frac{-10}{5} = 2\]

\textbf{The elasticity is still 2}, the same as in part (b). This is the same price rise, from \$20 to \$21, counted in cents. A 5\% rise is a 5\% rise whether we count in dollars or in cents, so the elasticity does not depend on the units.

\subsection{2. Two coffee shops}\label{two-coffee-shops}

\subsection{3. Flat and steep demand curves}\label{flat-and-steep-demand-curves}

\subsection{4. Which demand is more elastic?}\label{which-demand-is-more-elastic}

\subsection{5. Beach parking}\label{beach-parking}

\textbf{(a)} \[\begin{aligned}
\%\text{ change in price} &= \frac{25 - 20}{20} \times 100 = 25\% \\[0.3em]
\%\text{ change in cars} &= \frac{240 - 400}{400} \times 100 = -40\%
\end{aligned}\]
\[\varepsilon = -\frac{-40}{25} = 1.6\]

\textbf{Demand for the lot is elastic.}

\textbf{(b)} \[\varepsilon = -\frac{-5}{25} = 0.2\]

\textbf{Demand for beach visits is inelastic.}

\textbf{(c)} \textbf{Street parking is a close substitute for the lot}, so when the lot's price rises, many drivers switch. There is no close substitute for the beach itself, so most people still come. The same pattern showed up with Philadelphia's soda tax: stores just outside the city were close substitutes for stores inside it.

\subsection{6. Multiple choice}\label{multiple-choice}

\medskip

\textbf{Q1.} \textbf{(b) and (c).} On D, at \$4,000 the firm can sell 20 units, not 30. On D′, a price of \$4,000 corresponds to 60 units. At \$2,000 the firm sells 40 units on D and 80 units on D′, so 40 more. At 20 units the price is \$4,000 on D and \$8,000 on D′, a difference of \$4,000, not \$2,000. \textbf{(b) and (c)} are correct.

\medskip

\textbf{Q2.} \textbf{(a), (b) and (d).} \[\begin{aligned}
\%\text{ change in price} &= \frac{12 - 10}{10} \times 100 = 20\% \\[0.3em]
\%\text{ change in demand} &= \frac{15 - 20}{20} \times 100 = -25\%
\end{aligned}\]

So \(\varepsilon = -\dfrac{-25}{20} = 1.25\), which is greater than 1: demand is elastic, not inelastic. \textbf{(a), (b), and (d)} are correct.

\medskip

\textbf{Q3.} \textbf{(a) and (d).} At E both curves have the same price and quantity, and D₁ is steeper, so the same price rise loses fewer buyers on D₁: demand is less elastic there. Along one straight line the slope stays the same but the elasticity does not: at A the price is higher and the quantity lower than at C, so demand is more elastic at A. For the same reason, demand on D₂ is more elastic at E than at B. D₁ is less elastic, so its firm probably faces less competition, not more. \textbf{(a) and (d)} are correct.

\medskip

\textbf{Q4.} \textbf{(b).} \(\varepsilon = -\dfrac{-8}{4} = 2\). The minus sign in the formula makes the elasticity positive, and 2 is greater than 1, so demand is elastic. \textbf{(b)} is correct.

\medskip

\textbf{Q5.} \textbf{(a).} \[\begin{aligned}
\%\text{ change in price} &= \frac{25 - 20}{20} \times 100 = 25\% \\[0.3em]
\%\text{ change in demand} &= \frac{320 - 400}{400} \times 100 = -20\%
\end{aligned}\]

So \(\varepsilon = -\dfrac{-20}{25} = 0.8\), which is less than 1: demand is inelastic. Option (b) divides the wrong way round, and (d) divides the change in tickets by the change in price instead of using percentages. \textbf{(a)} is correct.

\section{Lecture 8: Steak, ramen, and recessions: Other elasticities}\label{lecture-8-steak-ramen-and-recessions-other-elasticities}

\subsection{1. Devon's month}\label{devons-month}

\textbf{(a)} Devon's income rises by
\[\frac{2{,}500 - 2{,}000}{2{,}000} \times 100 = 25\%\]

\begin{longtable}[]{@{}
  >{\raggedright\arraybackslash}p{(\linewidth - 8\tabcolsep) * \real{0.2000}}
  >{\raggedleft\arraybackslash}p{(\linewidth - 8\tabcolsep) * \real{0.2000}}
  >{\raggedleft\arraybackslash}p{(\linewidth - 8\tabcolsep) * \real{0.2000}}
  >{\raggedleft\arraybackslash}p{(\linewidth - 8\tabcolsep) * \real{0.2000}}
  >{\raggedleft\arraybackslash}p{(\linewidth - 8\tabcolsep) * \real{0.2000}}@{}}
\toprule\noalign{}
\begin{minipage}[b]{\linewidth}\raggedright
\end{minipage} & \begin{minipage}[b]{\linewidth}\raggedleft
At \$2,000
\end{minipage} & \begin{minipage}[b]{\linewidth}\raggedleft
At \$2,500
\end{minipage} & \begin{minipage}[b]{\linewidth}\raggedleft
\% change in demand
\end{minipage} & \begin{minipage}[b]{\linewidth}\raggedleft
Income elasticity
\end{minipage} \\
\midrule\noalign{}
\endhead
\bottomrule\noalign{}
\endlastfoot
Rides on the city bus & 40 & 30 & (30 − 40) / 40 × 100 = −25\% & −25 / 25 = −1 \\
Gallons of milk & 8 & 9 & (9 − 8) / 8 × 100 = 12.5\% & 12.5 / 25 = 0.5 \\
Rideshare trips & 2 & 3 & (3 − 2) / 2 × 100 = 50\% & 50 / 25 = 2 \\
\end{longtable}

There is no minus sign in the formula this time, so the sign of each answer is part of the answer.

\textbf{(b)} - \textbf{Bus rides are inferior for Devon.} The income elasticity is negative: he buys fewer of them as he gets richer.
- \textbf{Milk and rideshare trips are normal.} Both elasticities are positive.
- \textbf{Milk is a necessity}, at 0.5, which is between 0 and 1: Devon buys more milk as his income rises, but demand grows more slowly than income.
- \textbf{Rideshare trips are a luxury}, at 2, which is greater than 1: demand grows faster than income.

\textbf{(c)} The percentage change in demand is the elasticity times the percentage change in income.

\begin{itemize}
\tightlist
\item
  Bus rides: \(-1 \times (-10) = +10\%\), so he takes \textbf{about 10\% more} bus rides.
\item
  Milk: \(0.5 \times (-10) = -5\%\), so he buys \textbf{about 5\% less} milk.
\item
  Rideshare trips: \(2 \times (-10) = -20\%\), so he takes \textbf{about 20\% fewer} rideshare trips.
\end{itemize}

The luxury loses the most demand and the necessity the least. \textbf{Demand for the inferior good moves the other way and rises}, which is why some goods sell better in a recession.

\subsection{2. Carla's grocery bill}\label{carlas-grocery-bill}

\textbf{(a)} The price of ground beef rises by
\[\frac{5.00 - 4.00}{4.00} \times 100 = 25\%\]

\begin{longtable}[]{@{}
  >{\raggedright\arraybackslash}p{(\linewidth - 8\tabcolsep) * \real{0.2000}}
  >{\raggedleft\arraybackslash}p{(\linewidth - 8\tabcolsep) * \real{0.2000}}
  >{\raggedleft\arraybackslash}p{(\linewidth - 8\tabcolsep) * \real{0.2000}}
  >{\raggedleft\arraybackslash}p{(\linewidth - 8\tabcolsep) * \real{0.2000}}
  >{\raggedleft\arraybackslash}p{(\linewidth - 8\tabcolsep) * \real{0.2000}}@{}}
\toprule\noalign{}
\begin{minipage}[b]{\linewidth}\raggedright
\end{minipage} & \begin{minipage}[b]{\linewidth}\raggedleft
At \$4.00
\end{minipage} & \begin{minipage}[b]{\linewidth}\raggedleft
At \$5.00
\end{minipage} & \begin{minipage}[b]{\linewidth}\raggedleft
\% change in demand
\end{minipage} & \begin{minipage}[b]{\linewidth}\raggedleft
Cross-price elasticity
\end{minipage} \\
\midrule\noalign{}
\endhead
\bottomrule\noalign{}
\endlastfoot
Pounds of chicken & 10 & 15 & (15 − 10) / 10 × 100 = 50\% & 50 / 25 = 2 \\
Packs of hamburger buns & 8 & 6 & (6 − 8) / 8 × 100 = −25\% & −25 / 25 = −1 \\
Bottles of laundry detergent & 2 & 2 & 0\% & 0 / 25 = 0 \\
\end{longtable}

\textbf{(b)} - \textbf{Chicken and ground beef are substitutes.} The cross-price elasticity is positive: when beef gets more expensive, Carla buys chicken instead.
- \textbf{Hamburger buns and ground beef are complements.} The cross-price elasticity is negative: she buys the two together, so buying less beef means buying fewer buns.
- \textbf{Laundry detergent is neither.} Its cross-price elasticity is 0, so the price of beef tells us nothing about how much detergent she buys.

\subsection{3. Reading the elasticities}\label{reading-the-elasticities}

\textbf{(a)} - \textbf{Laundromat visits are an inferior good.} The income elasticity is negative, so households use the laundromat less as their income rises, presumably because they buy a washing machine.
- \textbf{Bottled water is a normal good and a necessity.} The elasticity is positive but between 0 and 1.
- \textbf{Cruise vacations are a normal good and a luxury.} The elasticity is greater than 1, so demand grows faster than income.
- \textbf{Tea and coffee are substitutes.} The cross-price elasticity is positive, so a rise in the price of coffee sends some buyers to tea.
- \textbf{Bagels and cream cheese are complements.} The cross-price elasticity is negative, so a rise in the price of cream cheese lowers demand for bagels.

\textbf{(b)} \textbf{Laundromat visits sell better.} Demand for an inferior good rises when income falls, so a 10\% fall in income raises laundromat visits by about 4\%.

\textbf{Cruise vacations lose the most.} They have the highest income elasticity, so the same 10\% fall in income cuts demand by about 22\%. Demand for bottled water falls by only about 3\%.

\subsection{4. Omar's taco truck}\label{omars-taco-truck}

\textbf{(a)} \[\%\text{ change in demand} = \frac{520 - 500}{500} \times 100 = 4\%\]
\[\varepsilon_{\text{cross}} = \frac{4}{10} = 0.4\]

\textbf{Tacos and burritos are substitutes}, because the elasticity is positive. They are not especially close substitutes: a 10\% price rise across the street moves only 4\% more tacos, so most of the burrito place's customers stay put.

\textbf{(b)} \[\varepsilon_{\text{income}} = \frac{4}{-8} = -0.5\]

\textbf{Omar's tacos are an inferior good for these customers.} The elasticity is negative: they buy more of his tacos when their incomes fall, probably because they are eating at cheaper places than before.

\subsection{5. Sorting goods}\label{sorting-goods}

\textbf{(a)} - \textbf{The plane ticket.} A vacation is a luxury people take when they can afford it, while the bus pass is what gets you to work at any income.
- \textbf{Dinner delivered from a restaurant.} People who come into money eat out more and cook cheap pasta less, so the pasta may even be inferior.
- \textbf{The weekend at a hotel.} Everyone needs somewhere to live, so rent grows more slowly than income, while hotel weekends grow faster.

\textbf{(b)} - \textbf{Butter and margarine: positive, substitutes.} A rise in the price of butter sends buyers to margarine.
- \textbf{Printers and ink cartridges: negative, complements.} They are used together, so a rise in the price of printers lowers demand for cartridges.
- \textbf{Gasoline and bus rides: positive, substitutes.} When gas gets expensive, some drivers take the bus instead.
- \textbf{Notebooks and lawn mowers: about zero, neither.} The price of one has no bearing on demand for the other.

\subsection{6. Multiple choice}\label{multiple-choice}

\medskip

\textbf{Q1.} \textbf{(b).} The income elasticity is \(\dfrac{-4}{10} = -0.4\). It is negative, so soup is an inferior good. Option (c) divides the wrong way round, and a luxury has an income elasticity greater than 1. \textbf{(b)} is correct.

\medskip

\textbf{Q2.} \textbf{(b).} Income rises by \(\frac{5{,}000 - 4{,}000}{4{,}000} \times 100 = 25\%\) and meals rise by \(\frac{6 - 4}{4} \times 100 = 50\%\), so the income elasticity is \(\dfrac{50}{25} = 2\). It is greater than 1, so restaurant meals are a luxury. \textbf{(b)} is correct.

\medskip

\textbf{Q3.} \textbf{(a), (b) and (d).} A luxury has an income elasticity greater than 1, not between 0 and 1, so (c) is wrong. The rest follow from the definitions: a negative income elasticity means demand falls as income rises, and so rises as income falls, and 0.4 is positive and below 1. \textbf{(a), (b), and (d)} are correct.

\medskip

\textbf{Q4.} \textbf{(b).} The cross-price elasticity is \(\dfrac{-16}{8} = -2\). It is negative, so the two goods are complements: the more expensive B gets, the less of A people buy. Option (c) divides the wrong way round. \textbf{(b)} is correct.

\medskip

\textbf{Q5.} \textbf{(c).} A negative cross-price elasticity means the goods are complements, bought together, so a rise in the price of one lowers demand for the other. Hot dogs and hot dog buns are the pair that go together. The other three pairs are substitutes, with positive cross-price elasticities. \textbf{(c)} is correct.

\medskip

\textbf{Q6.} \textbf{(a), (b) and (d).} The percentage change in demand is the elasticity times the percentage change in income, so recreation falls by about \(1.25 \times 10 = 12.5\%\) and food by about \(0.35 \times 10 = 3.5\%\). Recreation is above 1, which makes it a luxury, and food is below 1, which makes it a necessity, so (c) has them backwards. An inferior good has a negative income elasticity, so its demand rises when income falls. \textbf{(a), (b), and (d)} are correct.

\section{Lecture 9: Sweet spot: Choosing the price and quantity that maximize profit}\label{lecture-9-sweet-spot-choosing-the-price-and-quantity-that-maximize-profit}

\subsection{1. Sierra Bikes}\label{sierra-bikes}

\textbf{(a)} To sell 20 bikes, the demand curve says the price has to be \(P = 1{,}200 - 10 \times 20 = 1{,}000\). Revenue is \(R = P \times Q = 1{,}000 \times 20 = 20{,}000\). Cost is \(C(Q) = 2{,}000 + 200 \times 20 = 6{,}000\). So profit is \(20{,}000 - 6{,}000 = 14{,}000\).

Doing the same for the other quantities:

\begin{longtable}[]{@{}ccccc@{}}
\toprule\noalign{}
\(Q\) & \(P\) & \(R\) & \(C(Q)\) & Profit \\
\midrule\noalign{}
\endhead
\bottomrule\noalign{}
\endlastfoot
20 & 1,000 & 20,000 & 6,000 & 14,000 \\
30 & 900 & 27,000 & 8,000 & 19,000 \\
40 & 800 & 32,000 & 10,000 & 22,000 \\
50 & 700 & 35,000 & 12,000 & 23,000 \\
60 & 600 & 36,000 & 14,000 & 22,000 \\
\end{longtable}

\textbf{(b)} Profit is highest at \textbf{50 bikes a day}, where the price is \textbf{\$700} and profit is \textbf{\$23,000 a day}.

Revenue is higher still at 60 bikes, \$36,000 against \$35,000, but the extra 10 bikes add \$2,000 to cost and only \$1,000 to revenue, so profit falls.

\textbf{(c)} To sell 49 bikes, the price has to be \(1{,}200 - 10 \times 49 = 710\), so revenue is \(710 \times 49 = 34{,}790\). For 50 bikes it has to be \$700, so revenue is \(700 \times 50 = 35{,}000\). For 51 bikes it has to be \$690, so revenue is \(690 \times 51 = 35{,}190\).

The 50th bike adds \(35{,}000 - 34{,}790 = 210\) to revenue, more than its \$200 cost, so \textbf{it is worth selling}. The 51st adds only \(35{,}190 - 35{,}000 = 190\), less than it costs, so \textbf{it is not}. Sierra stops at 50 bikes, where marginal revenue drops below marginal cost, which matches part (b).

\textbf{(d)} \textbf{Neither changes.} The rent is paid whatever Sierra sells, so it does not change what one more bike adds to revenue or to cost: the 50th bike still adds \$210 against a cost of \$200, and the 51st only \$190. Profit falls by \$1,000 at every quantity, so it still peaks at \textbf{50 bikes at \$700}, and profit there is now \(23{,}000 - 1{,}000 =\) \textbf{\$22,000 a day}.

The rent does matter for a different decision: whether to operate at all. If Sierra could avoid the rent by closing, it compares what it earns from the bikes with the rent. At 50 bikes, revenue minus the cost of the bikes is \(35{,}000 - 200 \times 50 = 25{,}000\), so Sierra stays open as long as the rent is below \$25,000 a day.

\subsection{2. Marisol's cakes}\label{marisols-cakes}

\textbf{(a)} Revenue is \(P \times Q\), and cost is \(30 + 20Q\).

\begin{longtable}[]{@{}ccccc@{}}
\toprule\noalign{}
\(P\) & \(Q\) & \(R\) & \(C(Q)\) & Profit \\
\midrule\noalign{}
\endhead
\bottomrule\noalign{}
\endlastfoot
55 & 1 & 55 & 50 & 5 \\
50 & 2 & 100 & 70 & 30 \\
45 & 3 & 135 & 90 & 45 \\
40 & 4 & 160 & 110 & 50 \\
35 & 5 & 175 & 130 & 45 \\
30 & 6 & 180 & 150 & 30 \\
\end{longtable}

Profit is highest at \textbf{\$40 a cake}, where she sells \textbf{4 cakes} and makes \textbf{\$50 a day}.

\textbf{(b)}

\begin{longtable}[]{@{}ccccc@{}}
\toprule\noalign{}
Cake & Revenue goes from & to & \(\text{MR}\) & \(\text{MC}\) \\
\midrule\noalign{}
\endhead
\bottomrule\noalign{}
\endlastfoot
1st & 0 & 55 & 55 & 20 \\
2nd & 55 & 100 & 45 & 20 \\
3rd & 100 & 135 & 35 & 20 \\
4th & 135 & 160 & 25 & 20 \\
5th & 160 & 175 & 15 & 20 \\
6th & 175 & 180 & 5 & 20 \\
\end{longtable}

The first four cakes each add more to revenue than the \$20 they cost, so \textbf{the first four are worth baking}. The 5th cake adds only \$15, less than its cost. So Marisol stops at 4 cakes, the same answer as part (a).

Notice that marginal revenue is below the price: to sell the 4th cake she cuts the price from \$45 to \$40, so the 4th cake brings in \$40 but the other 3 cakes each sell for \$5 less, and \(40 - 15 = 25\).

\subsection{3. Sorting}\label{sorting}

\subsection{4. Multiple choice}\label{multiple-choice}

\medskip

\textbf{Q1.} \textbf{(a), (b) and (d).} The demand curve ties price and quantity together. To sell 100 smoothies the price has to be \(10 - 100/20 = 5\), and at \$6 buyers take \(20 \times (10 - 6) = 80\), not 120. So once the stand picks a quantity, the demand curve fixes the price, and the other way round. \textbf{(a), (b), and (d)} are correct.

\medskip

\textbf{Q2.} \textbf{(b).} Since \(Q\) is still 20, production costs remain the same, but revenue falls by \$200 on each car, or \$4,000 in total, so profit is \$340,000 − \$4,000 = \$336,000. Average cost is \(260{,}000 / 20 = \$13{,}000\); \$17,000 is the profit per car, \$340,000 / 20. At the lower price more than 20 buyers want a car, so the firm has no problem selling all 20. \textbf{(b)} is correct.

\medskip

\textbf{Q3.} \textbf{(a).} At a price above \$30,000 the firm can sell fewer than 20 cars, and it will not produce more cars than it can sell. Marginal cost is \$10,000 for every car, so it does not rise, and producing fewer cars lowers total cost. Profit falls, because 20 cars at \$30,000 was already the profit-maximizing choice. \textbf{(a)} is correct.

\medskip

\textbf{Q4.} \textbf{(a) and (c).} At 10 cars marginal revenue is \$30,000 and marginal cost is \$10,000, so one more car adds to profit. At 30 cars the extra car loses money, but the firm still makes a profit on the cars before it: profit at 30 cars is \$240,000, so it is not negative. Marginal revenue equals marginal cost at 20 cars, and the demand curve above that point gives the price, \$30,000. At 25 cars marginal revenue is zero, which is where revenue is highest, but every car past the 20th adds less to revenue than its \$10,000 cost, so profit is lower there. \textbf{(a) and (c)} are correct.

\medskip

\textbf{Q5.} \textbf{(c).} Up to 30 units, each extra unit adds more to revenue than the \$4 it costs, so profit rises. Past 30, marginal revenue is below \$4 and each extra unit lowers profit. Profit is highest where marginal revenue equals marginal cost, at 30 units. \textbf{(c)} is correct.

\medskip

\textbf{Q6.} \textbf{(b).} Revenue rises from \(30 \times 600 = 18{,}000\) to \(31 \times 590 = 18{,}290\), so marginal revenue is \$290. Put another way, the 31st unit brings in \$590, but the other 30 units each sell for \$10 less, a loss of \$300. \textbf{(b)} is correct.

\medskip

\textbf{Q7.} \textbf{(a).} While marginal revenue is above marginal cost, each extra unit adds to profit, so the firm expands until the two are equal. It cannot raise the price and hold quantity, because the demand curve ties the two together. \textbf{(a)} is correct.

\medskip

\textbf{Q8.} \textbf{(c) and (d).} The fixed cost is paid whatever the firm produces, so it does not change what one more unit adds to cost: marginal cost stays the same. Profit at every quantity falls by the same \$500, so the quantity with the highest profit, and the price that sells it, do not change. \textbf{(c) and (d)} are correct.

\section{Lecture 10: The surplus from a sale and who captures it}\label{lecture-10-the-surplus-from-a-sale-and-who-captures-it}

\subsection{1. Sierra Bikes}\label{sierra-bikes}

\textbf{(a)} The 10th buyer's willingness to pay is the height of the demand curve at 10 bikes: \(\text{WTP} = 1{,}200 - 10 \times 10 = 1{,}100\). Each bike costs Sierra \$200 to make, so the joint surplus is \(\text{WTP} - \text{MC} = 1{,}100 - 200 = 900\).

At a price of \$700, the buyer gets \(\text{WTP} - P = 1{,}100 - 700 =\) \textbf{\$400}, and Sierra gets \(P - \text{MC} = 700 - 200 =\) \textbf{\$500}. Together they add up to the joint surplus of \textbf{\$900}.

\textbf{(b)} \texttt{\{=latex\}\ \textbackslash{}begin\{center\}\ \textbackslash{}includegraphics{[}width=5.00in,alt=\{Sierra\ Bikes\textquotesingle{}\ demand\ curve,\ a\ straight\ line\ falling\ from\ 1,200\ dollars\ at\ zero\ bikes\ to\ zero\ at\ 120\ bikes,\ with\ marginal\ cost\ flat\ at\ 200\ dollars\ and\ the\ point\ E\ at\ 50\ bikes\ and\ 700\ dollars.\ The\ triangle\ between\ the\ demand\ curve\ and\ the\ price\ line\ at\ 700\ dollars,\ from\ zero\ to\ 50\ bikes,\ is\ shaded\ light\ blue\ and\ labelled\ consumer\ surplus,\ 12,500\ dollars.\ The\ rectangle\ between\ the\ price\ line\ and\ marginal\ cost\ at\ 200\ dollars,\ from\ zero\ to\ 50\ bikes,\ is\ shaded\ light\ orange\ and\ labelled\ producer\ surplus,\ 25,000\ dollars.\}{]}\{img/sierra-surplus.pdf\}\ \textbackslash{}end\{center\}}

Consumer surplus is the triangle between the demand curve and the price. The demand curve starts at \$1,200, so
\[\text{CS} = \tfrac{1}{2} \times 50 \times (1{,}200 - 700) = 12{,}500\]

Producer surplus is the rectangle between the price and marginal cost:
\[\text{PS} = (700 - 200) \times 50 = 25{,}000\]

Profit is producer surplus minus the fixed cost: \(25{,}000 - 2{,}000 =\) \textbf{\$23,000 a day}.

\subsection{2. Escape Hour}\label{escape-hour}

\textbf{(a)} Consumer surplus is the triangle between the demand curve (which starts at \$80) and the price:
\[\text{CS} = \tfrac{1}{2} \times 30 \times (80 - 50) = 450\]

Producer surplus is the rectangle between the price and marginal cost:
\[\text{PS} = (50 - 20) \times 30 = 900\]

\textbf{(b)} \textbf{Escape Hour}, with \$900 of the \$1,350. It is the only escape room in town, so it has market power: it can set a high price, and players who value the game highly will still pay it. One player cannot bargain for a better deal, because Escape Hour has many other customers.

\textbf{(c)} At \$40, the demand curve shows \textbf{40 tickets} sold.

\[\text{CS} = \tfrac{1}{2} \times 40 \times (80 - 40) = 800 \qquad \text{PS} = (40 - 20) \times 40 = 800\]

Consumer surplus \textbf{rises from \$450 to \$800}, and producer surplus \textbf{falls from \$900 to \$800}. The lower price moves surplus from Escape Hour to the players who were already buying, and the 10 extra players add surplus of their own.

\subsection{3. Multiple choice}\label{multiple-choice}

\medskip

\textbf{Q1.} \textbf{(b).} A buyer's willingness to pay is the height of the demand curve at their place in line: \(90 - 2 \times 15 = 60\). \textbf{(b)} is correct.

\medskip

\textbf{Q2.} \textbf{(b).} The buyer gets \(\text{WTP} - P = 120 - 80 = 40\). The store gets \(P - \text{MC} = 80 - 50 = 30\), and the joint surplus is \(120 - 50 = 70\). \textbf{(b)} is correct.

\medskip

\textbf{Q3.} \textbf{(a), (b) and (c).} Consumer surplus is the triangle \(\tfrac{1}{2} \times 20 \times (60 - 40) = 200\), and producer surplus is the rectangle \((40 - 20) \times 20 = 400\). Profit is producer surplus minus the fixed cost, \(400 - 150 = 250\). Price times quantity is revenue, not consumer surplus. \textbf{(a), (b), and (c)} are correct.

\medskip

\textbf{Q4.} \textbf{(c).} Producer surplus does not count the fixed cost, so profit is \(5{,}000 - 1{,}500 = 3{,}500\). \textbf{(c)} is correct.

\medskip

\textbf{Q5.} \textbf{(a), (b) and (d).} The joint surplus is \(\text{WTP} - \text{MC} = 1{,}100 - 200 = 900\), whatever the price. At \$700 the buyer gets \$400 and Sierra \$500; at \$650 the buyer gets \$450 and Sierra \$450. The price only moves surplus from one side to the other, so Sierra's gain falls. \textbf{(a), (b), and (d)} are correct.

\medskip

\textbf{Q6.} \textbf{(a).} The split of the surplus depends on bargaining power. The only ferry to an island can set a high price, and passengers who value the trip highly will still pay it. The other sellers each face many close competitors, so a high price would send their customers elsewhere. \textbf{(a)} is correct.

\section{Lecture 11: Deadweight loss and who captures the surplus}\label{lecture-11-deadweight-loss-and-who-captures-the-surplus}

\subsection{1. Sierra Bikes}\label{sierra-bikes}

\textbf{(a)} Demand meets marginal cost where \(1{,}200 - 10Q = 200\), at \textbf{\(Q = 100\) bikes}. Each of the first 100 bikes is worth more to its buyer than the \$200 it costs to make, but Sierra sells only 50.

The deadweight loss is the triangle between the demand curve and marginal cost from 50 to 100 bikes:
\[\text{DWL} = \tfrac{1}{2} \times (100 - 50) \times (700 - 200) = 12{,}500\]

The market could create \(\tfrac{1}{2} \times 100 \times (1{,}200 - 200) = 50{,}000\) of surplus a day, and at 50 bikes it creates \(12{,}500 + 25{,}000 = 37{,}500\).

\textbf{(b)} The 70th buyer's willingness to pay is the height of the demand curve at 70 bikes: \(1{,}200 - 10 \times 70 = 500\), less than \$700, so \textbf{this buyer does not buy}.

At \$350, the buyer gains \(500 - 350 =\) \textbf{\$150}, and Sierra gains \(350 - 200 =\) \textbf{\$150} on a bike it would not otherwise have sold. Everyone else still pays \$700, so nobody is worse off. \textbf{It is a Pareto improvement.}

\subsection{2. Escape Hour}\label{escape-hour}

\textbf{(a)} The 40th player's willingness to pay is the height of the demand curve at 40 tickets: \textbf{\$40}, less than \$50, so \textbf{this player does not buy}.

At \$30, the player gains \(40 - 30 =\) \textbf{\$10}, and Escape Hour gains \(30 - 20 =\) \textbf{\$10} on a ticket it would not otherwise have sold. Everyone else still pays \$50, so nobody is worse off. \textbf{It is a Pareto improvement.}

\textbf{(b)} The demand curve is above marginal cost up to \textbf{60 tickets}, but Escape Hour sells only 30. The deadweight loss is the triangle between the demand curve and marginal cost from 30 to 60 tickets:
\[\text{DWL} = \tfrac{1}{2} \times (60 - 30) \times (50 - 20) = 450\]

\subsection{3. Multiple choice}\label{multiple-choice}

\medskip

\textbf{Q1.} \textbf{(a) and (c).} An outcome is Pareto efficient when no change could make someone better off without making someone else worse off, so if a Pareto improvement exists, the outcome is not efficient. Giving the whole pizza to one person is efficient, because any slice for the others makes that person worse off, but few would call it fair. Moving \$10 helps the person who gets it and hurts the one who pays, so it is not a Pareto improvement. \textbf{(a) and (c)} are correct.

\medskip

\textbf{Q2.} \textbf{(a).} In (a) the new buyer gains, the firm earns more than the unit costs, and nobody else is affected. Raising the price hurts buyers, lowering it from the profit-maximizing price lowers the firm's profit, and dropping a buyer makes that buyer worse off. \textbf{(a)} is correct.

\medskip

\textbf{Q3.} \textbf{(c).} Past the profit-maximizing quantity, the extra revenue from one more unit is less than its cost, because the firm has to cut the price on all its units to sell it. There are still buyers who value the good at more than its marginal cost, which is why there is a deadweight loss. The fixed cost does not affect the choice of quantity. \textbf{(c)} is correct.

\medskip

\textbf{Q4.} \textbf{(a), (b) and (d).} Consumer surplus is \(\tfrac{1}{2} \times 40 \times (100 - 60) = 800\), and producer surplus is \((60 - 20) \times 40 = 1{,}600\). Demand meets marginal cost where \(100 - Q = 20\), at 80 units, so the deadweight loss is \(\tfrac{1}{2} \times (80 - 40) \times (60 - 20) = 800\), not \$1,600. \textbf{(a), (b), and (d)} are correct.

\medskip

\textbf{Q5.} \textbf{(a).} Producer surplus is revenue minus marginal costs, so it leaves out the fixed cost: profit is producer surplus minus the fixed cost. Deadweight loss is the loss of potential surplus to buyers and the firm together, not only the firm. At the profit-maximizing point some buyers are willing to pay more than a car costs to make but do not get one, so some gains from trade are missed. \textbf{(a)} is correct.

\medskip

\textbf{Q6.} \textbf{(a) and (d).} Price discrimination means charging different buyers different prices for the same good, based on how much they are willing to pay. A student discount and car-lot bargaining both do that. A large pizza is a different product from a small one, and an end-of-season sale gives every customer the same price. \textbf{(a) and (d)} are correct.

\section{Lecture 12: The markup rule and market power}\label{lecture-12-the-markup-rule-and-market-power}

\subsection{1. Canyon Kayaks}\label{canyon-kayaks}

\textbf{(a)} The profit margin is \(P - \text{MC} = 50 - 20 =\) \textbf{\$30} on each rental.

The markup is the profit margin as a share of the price:
\[\text{markup} = \frac{P - \text{MC}}{P} = \frac{30}{50} = 0.6\]
or \textbf{60\%} of the price.

\textbf{(b)} \[\varepsilon = -\frac{P}{Q} \times \frac{1}{\text{slope}} = \frac{50}{15} \times \frac{1}{2} = \frac{5}{3} \approx 1.67\]

\textbf{(c)} The markup rule says the markup equals \(1/\varepsilon\):
\[\frac{1}{\varepsilon} = \frac{3}{5} = 0.6\]
which is the \textbf{60\%} markup from part (a).

\textbf{(d)} Renters now have a close substitute, so more of them switch when Canyon Kayaks raises its price. Its demand becomes \textbf{more elastic}, and by the markup rule its \textbf{markup falls}.

\subsection{2. Phone cases}\label{phone-cases}

\textbf{(a)} The markup is \(1/\varepsilon = 1/4\), or \textbf{25\%}. Then
\[\frac{P - 15}{P} = \frac{1}{4}\]
so \(P - 15 = P/4\), which gives \(\tfrac{3}{4}P = 15\) and \textbf{\(P = 20\)}.

Check: \((20 - 15)/20 = 5/20 = 1/4\).

\textbf{(b)} The markup falls to \(1/6\):
\[\frac{P - 15}{P} = \frac{1}{6}\]
so \(P - 15 = P/6\), which gives \(\tfrac{5}{6}P = 15\) and \textbf{\(P = 18\)}. With a close substitute on the market, the company sets a lower price.

\subsection{3. Pixel Forge}\label{pixel-forge}

\textbf{(a)} Average cost is total cost divided by downloads, so each download carries its \$2 plus a share of the \$60 million:
\[\text{AC} = \frac{60{,}000{,}000}{Q} + 2\]

At 10 million downloads, \(\text{AC} = 6 + 2 =\) \textbf{\$8}. At 20 million, \(\text{AC} = 3 + 2 =\) \textbf{\$5}. At 60 million, \(\text{AC} = 1 + 2 =\) \textbf{\$3}.

It falls because the \$60 million is spread over more and more downloads.

\textbf{(b)} Profit is \((P - \text{AC}) \times Q\), so the firm breaks even where \(\text{AC} = P = 5\). From part (a), that is \textbf{20 million downloads}. With fewer, average cost is above \$5 and Pixel Forge loses money on every sale. With more, average cost is below \$5 and it makes a profit that grows with every extra download.

\textbf{(c)} - \textbf{One studio} sells 60 million downloads at an average cost of \$3, so its profit is \$(5 - 3) \times 60 = \$ \textbf{\$120 million}.
- \textbf{Three studios} sell 20 million each, at an average cost of \$5, so each \textbf{just breaks even}.
- \textbf{Six studios} sell 10 million each, at an average cost of \$8, so each \textbf{loses} \((5 - 8) \times 10 = -\$30\) \textbf{million}.

The market can support \textbf{at most three studios}. Each one has to pay the whole \$60 million but gets only a share of the buyers, so the more studios there are, the higher each one's average cost. When average cost falls with output, a market ends up with a few large firms.

\subsection{4. Multiple choice}\label{multiple-choice}

\medskip

\textbf{Q1.} \textbf{(a), (c) and (d).} Fewer substitutes, a patent, and stronger brand loyalty all make the firm's demand less elastic, so it can set a higher price without losing its buyers. Buyers who care mainly about price switch easily to rivals, which makes demand more elastic and lowers market power. \textbf{(a), (c), and (d)} are correct.

\medskip

\textbf{Q2.} \textbf{(b).} By the markup rule, the markup is \(1/\varepsilon = 1/5\), so \((50 - \text{MC})/50 = 1/5\). Then \(50 - \text{MC} = 10\), and \(\text{MC} = 40\). \textbf{(b)} is correct.

\medskip

\textbf{Q3.} \textbf{(c).} A natural monopoly comes from the cost structure: with large fixed costs, average cost keeps falling, so one firm serving the whole market has the lowest average cost. A patent and a strong brand are other sources of market power, and firms agreeing on prices form a cartel. \textbf{(c)} is correct.

\medskip

\textbf{Q4.} \textbf{(b).} There are many firms, entry is easy, and the products are differentiated, so each salon has some control over its price. \textbf{(b)} is correct. In perfect competition the products would be identical.

\medskip

\textbf{Q5.} \textbf{(b) and (c).} Close substitutes make demand more elastic, not less, so (a) is wrong and (b) is right. If the local shops close, the chain loses its rivals, its demand becomes less elastic, and it can raise prices, so (c) is right. Goods that are very different from the chain's have few close substitutes, so their demand is less elastic, not more. \textbf{(b) and (c)} are correct.

\medskip

\textbf{Q6.} \textbf{(b).} Loyal customers switch less readily when the price rises, so demand is less elastic. The markup rule, \((P - \text{MC})/P = 1/\varepsilon\), says a lower elasticity means a higher markup. \textbf{(b)} is correct.

\medskip

\textbf{Q7.} \textbf{(c).} A few firms, hard entry, and prices that depend on what rivals do are the marks of an oligopoly. Monopolistic competition has many firms and easy entry, and a monopoly has one seller. \textbf{(c)} is correct.

\medskip

\textbf{Q8.} \textbf{(b).} In a natural monopoly one firm supplies the market at a lower average cost than several could, so splitting it up would raise costs. Left alone, a monopoly sets a high markup. A price equal to marginal cost would not cover the fixed cost of the pipes, so the utility could not survive. Regulating the price of a single utility, or public ownership, is the usual answer. \textbf{(b)} is correct.

\medskip

\textbf{Q9.} \textbf{(a), (b) and (d).} The training cost is a fixed cost, spread over every question answered, so average cost falls with use and stays above the few cents of marginal cost. Charging only marginal cost would leave nothing to pay for the training, so the company would lose the billions it spent. A newcomer with few users spreads its own training cost over fewer questions, so its average cost is higher and it cannot survive at a price the established company can. This is why markets like this tend to end up with a few large firms. \textbf{(a), (b), and (d)} are correct.

\medskip

\textbf{Q10.} \textbf{(a), (b) and (c).} Buying a rival, forming a cartel, and blocking rivals' access to buyers all reduce competition, and competition policy watches for each of them. Passing a cost saving on to buyers in a lower price does not reduce competition. \textbf{(a), (b), and (c)} are correct.

\section{Lecture 13: Setting prices with few rivals}\label{lecture-13-setting-prices-with-few-rivals}

\subsection{Multiple choice}\label{multiple-choice}

\medskip

\textbf{Q1.} \textbf{(b).} With only three sellers, a price cut by one takes customers from the other two, who then have to decide whether to follow. With hundreds or thousands of sellers, any one firm is too small for its price to matter to the others, and a single utility has no rivals at all. \textbf{(b)} is correct.

\medskip

\textbf{Q2.} \textbf{(c).} A firm with many small rivals can treat its demand curve as fixed, because its own price barely affects anyone else. With one or two rivals, a rival's price cut pulls customers away and shifts the firm's demand curve, so the firm has to think about what the rival will do before setting its own price. \textbf{(c)} is correct.

\medskip

\textbf{Q3.} \textbf{(a), (b) and (c).} Loyal customers make each shop's demand less elastic, so a high price loses few of them, and both shops can charge more. Price-sensitive customers move to whichever shop is cheaper, which pushes both prices down. And because a price cut by one shop takes customers from the other, each shop's best price depends on what it expects its rival to do. \textbf{(a), (b), and (c)} are correct.

\medskip

\textbf{Q4.} \textbf{(a).} When the shops differ more, fewer customers switch when one raises its price, so each shop's demand is less elastic. Less competition between them means both can charge more. \textbf{(a)} is correct.

\medskip

\textbf{Q5.} \textbf{(c).} Both airlines earn less than they would at high fares, while passengers pay less. Competition between the two works in the passengers' favor. \textbf{(c)} is correct.

\medskip

\textbf{Q6.} \textbf{(a), (b) and (c).} A cartel is a group of firms that agree to act like one monopoly, setting a high price together instead of competing. That is why cartels are illegal in most countries: the higher fares come out of passengers' pockets. \textbf{(a), (b), and (c)} are correct.

\end{document}
